\hmsection{Conclusion}{}{}

This project developed an autonomous robotic system for
detecting, classifying, and sorting coloured balls using
computer vision and a four-degree-of-freedom arm with a fifth
motor for the gripper. The system integrates
a deterministic HSV-based detection pipeline running on a
Raspberry Pi~5, a closed-form geometric inverse kinematics
solver, an adaptive current-sensing grip verification mechanism,
and a dedicated firmware layer on the OpenRB-150 microcontroller.

The minimum viable product defined in Section~\ref{sec:mvp} was achieved
under the tested laboratory conditions. In the final 20-ball
validation session, the system sorted \textbf{20/20 balls
correctly (100\% observed net sorting success)} with \textbf{0
unrecovered failures} (Section~\ref{sec:mvp-verification}). Both colour classes were
completed successfully: 10/10 blue balls and 10/10 red balls were
picked and placed in the correct bins. This net sorting result is
distinct from grip-attempt behaviour: one additional
grip-verification rejection occurred before successful retry. Grip
verification correctly rejected an unstable grasp, and retry logic
recovered the cycle on the next scan, so the event did not reduce
the final net sorting result.

Several design decisions proved effective. The closed-form
geometric IK solver provided deterministic, sub-millisecond
solutions without the convergence issues of numerical methods.
The adaptive grip algorithm, which exploits the Dynamixel
motor's current and load registers as an impedance-based force
sensor, reliably distinguished between a gripped ball and an
empty close without requiring external sensor hardware. The
touch calibration procedure superseded the polynomial sag model
by providing per-point surface heights that account for the
arm's real-world droop at each calibrated position.

The principal limitations identified during the project are the
remaining movement-related thermal load, the absence of
mid-approach visual correction due to the wrist-mounted camera
design, and the four-degree-of-freedom constraint that restricts
the gripper to a fixed downward orientation. The revised scan
pose mitigated idle scan-pose heating, but movement-related
heating remains an operational limitation and risk. The most
impactful improvement identified is to reduce the mass of the
wrist-to-claw assembly, which would directly lower the
gravitational torque on the elbow motor and reduce the arm's
overall sag (Section~\ref{sec:improvements}).
